% GENERATED FILE. Edit canonical lessons, then run npm run generate:books.
% canonical-source-hash: fnv1a64:d95ab84e18de95a3
% canonical-lessons: TA-W283-a2-connected-composition

\chapter{A connected A2 Tamil message}
\label{ch:ta-a-connected-a2-tamil-message}
\hlchaptermodality{283}

\begin{chapteropening}
\textbf{By the end of this chapter:} \emph{I can plan and write a connected 70–90 word Tamil message in Tamil script for a named reader and practical purpose.}

Complete chapter 283 by writing a model-free 70–90 word connected Tamil message for a named reader and practical purpose.
\end{chapteropening}

\section[(connected A2 Tamil writing)]{A connected Tamil message about your familiar day}

\label{lesson:TA-W283-a2-connected-composition}



\begin{quote}
A connected message is not a pile of vocabulary and not one long sentence. It is a short route for one reader. Before writing, draw eight empty boxes and give them these jobs:

\begin{enumerate}
\raggedright
  \item reader and greeting;
  \item reason for writing;
  \item where or when your day begins;
  \item two familiar actions;
  \item what happens afterwards and a second time anchor;
  \item one contrast or reason;
  \item one direct question about tomorrow;
  \item a close.
\end{enumerate}

Write only one or two English cue words in each box. Do not draft Tamil yet.
\end{quote}

\subsection*{Writing — connected composition}
There is no Tamil model, sentence bank, word bank, romanization, or copyable answer on this page.

\textbf{Reader:} your friend Meena. \textbf{Purpose:} help Meena understand your familiar day and choose a good time to visit you tomorrow.

Write \textbf{70–90 words in Tamil script}. Use every planning box. Include at least four familiar actions, two time anchors, one sequence link, one contrast or reason, and one direct question for Meena. Make every detail help with tomorrow's visit. The instructions give content, not sentences: choose only words and forms this book has taught, and connect them yourself.

This first connected attempt is untimed. Put away dictionaries, grammar references, translators, generative assistants, and spell-checkers. Do not copy from an earlier lesson. Ordinary legible Tamil handwriting is enough.

When the whole message exists, count its words once. If it is short, add one useful taught detail to a planning box. If it is long, remove a detail that does not help Meena. Do not replace the message with a memorised model.

\begin{tcolorbox}[breakable,colback=teal!4,colframe=teal!35!black,title={Before you move on}]
Review the same attempt in four separate passes:

\begin{enumerate}
\raggedright
  \item \textbf{Task fulfilment and relevant content} — Meena, the familiar day, tomorrow's visit, and every requested content point are present.
  \item \textbf{Organisation and cohesion} — the eight jobs form a route rather than a list of disconnected sentences.
  \item \textbf{Tamil vocabulary and grammar} — every form and connector comes from the taught standard-colloquial Tamil track; a simpler familiar sentence is better than an invented structure.
  \item \textbf{Tamil-script control} — letter shapes, vowel signs, conjuncts, word boundaries, and punctuation remain readable.
\end{enumerate}

Mark one strength and one next step. Repair only the next step; keep the original attempt so the next 70–90-word message can show real growth.

This is one untimed curriculum practice attempt, \textbf{not the complete 30-minute A2 writing paper} and not evidence that either timed mock has been passed. Closing the machine-measured connected-composition stage does not prove A2 attainment or readiness, learner performance, independent review, standard setting, mock passage, calibration, or human validation.
\end{tcolorbox}
